\section{Discussion}
\label{sec:discussion}

The formative and evaluation results together reposition skill repair from
editing a natural-language artifact to governing a change in recurring agent
behavior. Participants could often recognize an unacceptable outcome before
they could state the reusable policy that should replace it. The system study
suggests that resolving this articulation gap requires more than an improved
text editor or a larger set of automatically generated tests. Workflow owners
benefited from keeping their intended behavioral delta, the candidate's
artifact delta, and its enacted behavioral delta available as related but
non-equivalent objects of review.

\subsection{The Unit of Review Is a Behavioral Change, Not a Textual Patch}

Code review works partly because the reviewed artifact has a comparatively
direct relationship to later execution. Natural-language skills weaken that
relationship: a small edit may interact with the task context, other
instructions, tool results, and model variation, while a substantial textual
change may leave observable behavior untouched. Existing prompt and agent
evaluation systems make important parts of this uncertainty testable
\cite{arawjo2024chainforge,kim2024evallm,sharma2025promptpex}, but their
comparison becomes most decisive once the desired behavior is sufficiently
specified.

Behavioral patch review instead treats the textual diff as one representation
within a larger review. The scope states the change the owner is prepared to
authorize; the artifact diff shows how that change was encoded; and matched
execution provides bounded evidence about what the candidate actually does.
The improvement in held-out prediction suggests that this organization helped
participants form a more operational understanding of the revision than the
same capabilities presented through chat. This is not evidence that textual
diffs are unhelpful. Rather, familiar artifact review becomes more useful when
it remains connected to the cases and executions that give the edited text
behavioral meaning.

This framing also changes what counts as a successful repair. Reproducing the
desired outcome in the incident is necessary but insufficient. A candidate is
reviewable only relative to behavior the owner intended to change, behavior
they intended to preserve, and boundaries for which no rule has yet been
authorized. The relevant unit is therefore a \emph{relative behavioral
change}, not an isolated outcome or a globally optimal skill.

\subsection{Repair Scope Is a Provisional Behavioral Specification}

The formative study showed that repair intent was constructed comparatively:
additional cases did not merely test a pre-existing policy, but exposed
preservation commitments and trade-offs that participants had not articulated
from the incident alone. The controlled study extends this observation. The
workspace produced stronger scope adherence and regression detection even
though both conditions had access to the same cases and agent operations. The
result suggests that cases can function as an interactive specification when
their relationships and normative status remain visible across revision and
execution.

This specification is deliberately provisional. \textsc{Change} and
\textsc{Preserve} express commitments that can be evaluated against observable
behavior, while \textsc{Unresolved} records the absence of an authorized rule.
Treating unresolved behavior as a first-class state avoids two common
collapses: allowing the candidate generator to silently settle a domain
trade-off, or forcing the owner to invent a complete policy before any
evidence is available. New executions may legitimately alter the scope, but
preserving pre-reveal judgments distinguishes learning from retrofitting the
specification to a candidate that has already been observed.

This account complements policy-authoring and test-generation systems
\cite{lam2025policymaps,lee2026dataprompt,koh2026contrastive}. A living test set
can evaluate a stable contract, whereas a behavioral scope also records how
the contract is still being negotiated. For reusable agents, future authoring
tools should therefore preserve the history and provenance of case-level
commitments, not only accumulate examples labeled successful or unsuccessful.

\subsection{Evidence Can Support Authority Without Replacing It}

Participants in the formative study distinguished work they were willing to
delegate from decisions they wanted to retain. They generally accepted AI
support for retrieving cases, running comparisons, and drafting skill text,
while treating acceptable policy and release as owner decisions. The user
study's decision results suggest that a structured evidence record can make
this division of labor operational: participants made decisions that more
closely reflected mismatches and evidence gaps without losing perceived
control.

The distinction is important because additional evidence is not equivalent to
additional authority. A source intervention may reveal sensitivity to one
instruction but cannot determine whether the resulting policy is acceptable.
Matched executions may show that a candidate crossed an unresolved boundary
but cannot resolve the underlying trade-off. Even a candidate that satisfies
every reviewed case may behave differently outside that scope. \systemname{}
therefore automates evidence and implementation work while preserving explicit
moments at which the owner must interpret, revise, or decline the proposed
change.

This design differs from both full automation and manual oversight of every
agent step. Full automation risks turning incomplete evidence into an implicit
policy, while continuous manual supervision undermines the value of reusable
agents. Behavioral patch review concentrates owner attention at consequential
boundaries: defining the intended change, inspecting conflicts between intent
and execution, and authorizing persistent reuse. In this sense, the system
supports governance through selective, evidence-grounded intervention rather
than permanent human involvement in each workflow.

\subsection{Implications for Reusable Agent Maintenance}

Our findings suggest four implications for systems that maintain reusable
agent artifacts. First, repair interfaces should preserve intent separately
from implementation. A candidate author can explain how text is meant to
address a commitment, but that rationale should not overwrite the owner's
wording or count as behavioral evidence. Second, systems should compare the
complete revised artifact under matched conditions. Local source or block
interventions can guide attention, but interactions among instructions make
whole-candidate execution the relevant basis for release review.

Third, evidence should retain its scope and provenance. Cases, artifact
versions, model and tool configurations, and exposure to candidate authors all
shape what a comparison can support. A compact interface may hide these fields
by default, but it should not detach summaries from reconstructable records.
Finally, release should create reusable review material rather than end the
process. Accepted Change and Preserve commitments can become regression cases,
while Unresolved cases remain visible boundaries for later evidence and policy
work.

Although we instantiate these implications through agent skills, the model may
apply to other natural-language artifacts that govern recurring behavior,
including system prompts, organizational policies, automation routines, and
agent templates. The key requirement is not a particular file format, but a
versionable control artifact whose effects can be observed across multiple
executions and judged by a person with authority over reuse. Determining how
behavioral patch review changes when authority is distributed across teams or
institutions remains an important direction for future work.
